\documentclass[10pt,a4paper]{amsart}
\usepackage[margin=19mm]{geometry}
\usepackage{amsmath,amssymb,mathtools}
\usepackage[scr=rsfs,cal=euler]{mathalfa}
\usepackage{xfrac}
\usepackage[unicode,hidelinks,bookmarks=false]{hyperref}
\newcommand{\ie}{i.e.\ }
\newcommand{\cf}{cf.\ }
\newcommand{\resp}{respectively\ }
\newcommand{\Subsection}{\subsection}
\providecommand{\nobreakdash}{\nobreak\hbox{-}}
\setlength{\emergencystretch}{2em}
\multlinegap0pt
\usepackage{bm}
\usepackage{bbm}
\usepackage{dsfont}
\usepackage{xspace}
\usepackage{cancel}
\usepackage{mathtools}
\usepackage{amsthm}
\makeatletter
\@ifundefined{theorem}{\newtheorem{theorem}{Theorem}}{}
\@ifundefined{corollary}{\newtheorem{corollary}[theorem]{Corollary}}{}
\makeatother
\makeatletter
\expandafter\ifx\csname proof\endcsname\relax
\newenvironment{proof}[1][Proof]{\noindent\textbf{#1.} }{\ \rule{0.5em}{0.5em}}
\fi
\makeatother
\title[Compactness of products of block Hankel and Toeplitz operato]{Compactness of products of block Hankel and Toeplitz operators}
\author{Caixing Gu, Meng Li, Pan Ma}
\date{Source: arXiv:2511.12542v1 (CC BY 4.0)}
\begin{document}
\maketitle

\noindent\emph{Source and licence.} Compactness of products of block Hankel and Toeplitz operators. Version arXiv:2511.12542v1 (CC BY 4.0), distributed under the Creative Commons Attribution 4.0 International licence, as recorded in the arXiv licence metadata for this exact version and shown on the abstract page https://arxiv.org/abs/2511.12542v1. Full rights notice: licenses/arxiv-2511.12542-CC-BY-4.0.txt.\par\smallskip

\noindent\textit{Editorial scope.} Counted unit: the Corollary whose source label is \texttt{cor2} in the article (source0.tex lines 2753--2774, source label \texttt{cor2}), delivered together with the article's own complete original proof of it (source0.tex lines 2776--3183, one \texttt{proof} environment of 408 lines). No mathematics is written, repaired or re-derived: statement and proof are the article's own text line for line. Required context retained uncounted: maina (lines 267--280). Every definition, display and label of those lines is retained unchanged. Omitted from this package and not redistributed: 1--266, 281--2752, 2775--2775, 3184--3399. Nothing inside a retained excerpt is omitted or altered: every retained body is delivered exactly as the article's lines are written, so no omission receipt of any kind accompanies this package. Citations: the retained excerpts carry no citation command at all, so no bibliography entry of the article is transcribed and no reference is redistributed. Reference adaptation: none was needed and none was made; every reference inside the delivered unit resolves inside it, so no unresolved reference or citation is produced. Printed slips kept literal: no printed word, symbol, hypothesis or number of the counted statement or of its proof was altered. Layout and class: the article's own class is \texttt{article} with packages including amsmath, amsfonts, amssymb, graphicx, color, hyperref; the wrapper is the lane's pinned \texttt{amsart} editorial wrapper at 10pt on A4, which restates the theorem-like environments the retained text needs and adds the article's own macro definitions that the retained text needs (none), with extra wrapper packages (bm, bbm, dsfont, xspace, cancel, mathtools). The delivered numbering is the unit's own. Assets: no retained span contains a figure environment, a graphics-inclusion command, a PDF-page inclusion, a table or a TikZ picture; no image file is copied into this package, the package declares no asset dependency and no image file is redistributed with it. Licence: version arXiv:2511.12542v1 of this article is distributed under the Creative Commons Attribution 4.0 International licence, as recorded in the arXiv licence metadata for this exact version and shown on the abstract page https://arxiv.org/abs/2511.12542v1. A full rights notice is delivered at licenses/arxiv-2511.12542-CC-BY-4.0.txt. Characters: every retained excerpt is pure ASCII, so the delivered engine, XeLaTeX with its default Latin Modern text font, reports no missing character.
\begin{theorem}
\label{maina}Assume $E=\mathbb{C}^{n}$ and $\Phi,\Psi\in L_{B(E)}^{\infty}.$
Then $H_{\Phi}T_{\Psi}$ is compact on the vector-valued Hardy space $H_{E}%
^{2}$ if and only if
\[
\bigcap_{A\in\left(  M_{n}\right)  _{2^{2n}}}H^{\infty}\left[  \Phi
(I-A),A\Psi,\Phi A\Psi\right]  \subset H_{n\times n}^{\infty}+C_{n\times n},
\]
where \[
\bigcap_{A \in (M_{n})_{2^{2n}}} 
   H^{\infty}\!\bigl[\Phi(I-A),\, A\Psi,\, \Phi A\Psi\bigr]
\]denotes the intersection, over all $A \in (M_{n})_{2^{2n}}$, 
of the Douglas algebras generated by $H^{\infty}$ together with $\Phi(I-A)$, $A\Psi$ and $\Phi A\Psi$.
\end{theorem}

\begin{corollary}
\label{cor2} Let $\varphi_{i},\psi_{i}\in L^{\infty}$ for $i=1,2.$ Then
$H_{\varphi_{1}}T_{\psi_{1}}+H_{\varphi_{2}}T_{\psi_{2}}$ is compact if and
only if
\begin{align}
&  H^{\infty}\left[  \varphi_{1},\varphi_{2}\right]  \cap H^{\infty}\left[
\psi_{1},\psi_{2},\varphi_{1}\psi_{1}+\varphi_{2}\psi_{2}\right]
\label{line1}\\
&  \cap H^{\infty}\left[  \varphi_{2},\psi_{1},\varphi_{1}\psi_{1}\right]
\label{line-1}\\
&  \cap H^{\infty}\left[  \varphi_{1},\psi_{2},\varphi_{2}\psi_{2}\right]
\label{line0}\\
&  \cap_{\lambda\neq0}H^{\infty}\left[  \varphi_{2}-\lambda\varphi_{1}%
,\psi_{1}+\lambda\psi_{2},\varphi_{1}\left(  \psi_{1}+\lambda\psi_{2}\right)
\right] \label{line2}\\
&  \cap_{\lambda\neq0}H^{\infty}\left[  \varphi_{2}-\lambda\varphi_{1}%
,\psi_{1}+\lambda\psi_{2},\varphi_{2}\left(  \psi_{1}+\lambda\psi_{2}\right)
\right] \label{line3}\\
&  \subset H^{\infty}+C.\nonumber
\end{align}

\end{corollary}

\begin{proof}
Write
\begin{align}
\Phi=\left[
\begin{array}
[c]{cc}%
\varphi_{1} & \varphi_{2}\\
0 & 0
\end{array}
\right]  ,\Psi=\left[
\begin{array}
[c]{cc}%
\psi_{1} & 0\\
\psi_{2} & 0
\end{array}
\right]  .\nonumber
\end{align}
By Theorem \ref{maina}, $H_{\varphi_{1}}T_{\psi_{1}}+H_{\varphi_{2}}%
T_{\psi_{2}}$ is compact if and only if
\[
\bigcap_{A\in\left(  M_{2}\right)  _{2^{4}}}H^{\infty}\left[  \Phi
(I-A),A\Psi,\Phi A\Psi\right]  \subset H^{\infty}_{2\times2}+C_{2\times2}.
\]


\textbf{{Case A:}} In the case $rankA=2$ or $rank(I-A)=2,$ we have%
\[
H^{\infty}\left[  \Phi(I-A),A\Psi,\Phi A\Psi\right]  = H^{\infty}\left[
\Psi,\Phi\Psi\right]  \text{ or }H^{\infty}\left[  \Phi(I-A),A\Psi,\Phi
A\Psi\right]  =H^{\infty}\left[  \Phi\right]  ,
\]
i.e.
\[
H^{\infty}\left[  \Phi(I-A),A\Psi,\Phi A\Psi\right]  =H^{\infty}\left[
\left[
\begin{array}
[c]{cc}%
\psi_{1} & 0\\
\psi_{2} & 0
\end{array}
\right]  , \left[
\begin{array}
[c]{cc}%
\varphi_{1}\psi_{1}+\varphi_{2}\psi_{2} & 0\\
0 & 0
\end{array}
\right]  \right]  \quad\text{or}\quad H^{\infty} \left[
\begin{array}
[c]{cc}%
\varphi_{1} & \varphi_{2}\\
0 & 0
\end{array}
\right]  ,
\]
which is the same Douglas algebras as (\ref{line1}).

\textbf{{Case B:}} Next, it suffices to discuss $A\in\left(  M_{2}\right)
_{16}$ such that $rankA=rank(I-A)=1.$ In this situation, there are three
cases:%
\begin{align}
A  &  =\left[
\begin{array}
[c]{cc}%
1 & \lambda\\
0 & 0
\end{array}
\right]  ,\label{line1a}\\
A  &  =\left[
\begin{array}
[c]{cc}%
0 & 0\\
\lambda & 1
\end{array}
\right]  ,\label{line2a}\\
A  &  =\left[
\begin{array}
[c]{cc}%
1-\lambda t & \lambda\left(  1-\lambda t\right) \\
t & \lambda t
\end{array}
\right]  ,\left\vert \lambda\right\vert \leq16. \label{line3a}%
\end{align}


\textbf{{Case 1:}} If $\lambda=0$ in the matrix (\ref{line1a}), we have
\[
H^{\infty}\left[  \Phi(I-A),A\Psi,\Phi A\Psi\right]  =H^{\infty} \left[
\left[
\begin{array}
[c]{cc}%
0 & \varphi_{2}\\
0 & 0
\end{array}
\right]  ,\left[
\begin{array}
[c]{cc}%
\psi_{1} & 0\\
0 & 0
\end{array}
\right]  ,\left[
\begin{array}
[c]{cc}%
\varphi_{1}\psi_{1} & 0\\
0 & 0
\end{array}
\right]  \right]  ,
\]
which is the same Douglas algebras as (\ref{line-1}).

If $\lambda\neq0$ in the matrix (\ref{line1a}),%
\begin{align}
\Phi(I-A)=\left[
\begin{array}
[c]{cc}%
\varphi_{1} & \varphi_{2}\\
0 & 0
\end{array}
\right]  \left[
\begin{array}
[c]{cc}%
0 & -\lambda\\
0 & 1
\end{array}
\right]  =\left[
\begin{array}
[c]{cc}%
0 & \varphi_{2}-\lambda\varphi_{1}\\
0 & 0
\end{array}
\right]  ,\nonumber
\end{align}
\begin{align}
A\Psi= \left[
\begin{array}
[c]{cc}%
1 & \lambda\\
0 & 0
\end{array}
\right]  \left[
\begin{array}
[c]{cc}%
\psi_{1} & 0\\
\psi_{2} & 0
\end{array}
\right]  =\left[
\begin{array}
[c]{cc}%
\psi_{1}+\lambda\psi_{2} & 0\\
0 & 0
\end{array}
\right]  ,\nonumber
\end{align}
%

\begin{align}
\Phi A\Psi=\left[
\begin{array}
[c]{cc}%
\varphi_{1} & \varphi_{2}\\
0 & 0
\end{array}
\right]  \left[
\begin{array}
[c]{cc}%
\psi_{1}+\lambda\psi_{2} & 0\\
0 & 0
\end{array}
\right]  =\left[
\begin{array}
[c]{cc}%
\varphi_{1}\left(  \psi_{1}+\lambda\psi_{2}\right)  & 0\\
0 & 0
\end{array}
\right]  .\nonumber
\end{align}
Therefore, we have
\[
H^{\infty}\left[  \Phi(I-A),A\Psi,\Phi A\Psi\right]  =H^{\infty} \left[
\left[
\begin{array}
[c]{cc}%
0 & \varphi_{2}-\lambda\varphi_{1}\\
0 & 0
\end{array}
\right]  ,\left[
\begin{array}
[c]{cc}%
\psi_{1}+\lambda\psi_{2} 0 & \\
0 & 0
\end{array}
\right]  ,\left[
\begin{array}
[c]{cc}%
\varphi_{1}\left(  \psi_{1}+\lambda\psi_{2}\right)  & 0\\
0 & 0
\end{array}
\right]  \right]  ,
\]
which is the same Douglas algebras as (\ref{line2}).

\textbf{{Case 2:}} If $\lambda=0$ in the matrix (\ref{line2a}), we have
\[
H^{\infty}\left[  \Phi(I-A),A\Psi,\Phi A\Psi\right]  =H^{\infty} \left[
\left[
\begin{array}
[c]{cc}%
\varphi_{1} & 0\\
0 & 0
\end{array}
\right]  ,\left[
\begin{array}
[c]{cc}%
0 & 0\\
\psi_{2} & 0
\end{array}
\right]  ,\left[
\begin{array}
[c]{cc}%
\varphi_{2}\psi_{2} & 0\\
0 & 0
\end{array}
\right]  \right]  ,
\]
which is the same Douglas algebras as (\ref{line0}).

If $\lambda\neq0$ in the matrix (\ref{line2a}),
\begin{align}
\Phi(I-A)=\left[
\begin{array}
[c]{cc}%
\varphi_{1} & \varphi_{2}\\
0 & 0
\end{array}
\right]  \left[
\begin{array}
[c]{cc}%
1 & 0\\
-\lambda\  & 0
\end{array}
\right]  =\left[
\begin{array}
[c]{cc}%
-\lambda\left(  \varphi_{2}-\frac{1}{\lambda}\varphi_{1}\right)  & 0\\
0 & 0
\end{array}
\right]  ,\nonumber
\end{align}
%

\begin{align}
A\Psi= \left[
\begin{array}
[c]{cc}%
0 & 0\\
\lambda & 1
\end{array}
\right]  \left[
\begin{array}
[c]{cc}%
\psi_{1} & 0\\
\psi_{2} & 0
\end{array}
\right]  =\left[
\begin{array}
[c]{cc}%
0 & 0\\
\lambda\left(  \psi_{1}+\frac{1}{\lambda}\psi_{2}\right)  & 0
\end{array}
\right]  ,\nonumber
\end{align}
%

\begin{align}
\Phi A\Psi &  =\left[
\begin{array}
[c]{cc}%
\varphi_{1} & \varphi_{2}\\
0 & 0
\end{array}
\right]  \left[
\begin{array}
[c]{cc}%
0 & 0\\
\lambda\left(  \psi_{1}+\frac{1}{\lambda}\psi_{2}\right)  & 0
\end{array}
\right]  =\left[
\begin{array}
[c]{cc}%
\lambda\varphi_{2} \left(  \psi_{1}+\frac{1}{\lambda}\psi_{2}\right)  & 0\\
0 & 0
\end{array}
\right]  .\nonumber
\end{align}
Therefore, we have \begin{flalign}
&H^{\infty}\left[  \Phi(I-A),A\Psi,\Phi A\Psi\right] \nonumber
\\ \nonumber&=H^{\infty} \left[ \left[
%
\begin{array}
[c]{cc}%
-\lambda\left( \varphi_{2}-\frac{1}{\lambda}\varphi_{1}\right)  & 0  \\
0 & 0
\end{array}
\right] ,\left[
\begin{array}
[c]{cc}%
0 & 0  \\
\lambda\left( \psi_{1}+\frac{1}{\lambda}\psi_{2}\right)  & 0
\end{array}
\right] ,\left[
\begin{array}
[c]{cc}%
\lambda\varphi_{2} \left( \psi_{1}+\frac{1}{\lambda}\psi_{2}\right)  & 0  \\
0 & 0
\end{array}
\right] \right] ,
\end{flalign}
which is the same Douglas algebras as (\ref{line3}).

\textbf{{Case 3:}} We will show that the matrix (\ref{line3a}) does not yield
new Douglas algebras.

\textbf{{Case 3.1:}} When $t=0$, the matrix $A$ is (\ref{line1a}). This case
is reduced to \textbf{{Case 1}}

\textbf{{Case 3.2:}} When $1-\lambda t=0$, the matrix $A$ is $\left[
\begin{array}
[c]{cc}%
0 & 0\\
\frac{1}{\lambda} & 1
\end{array}
\right]  $, and we get the same Douglas algebras as the case (\ref{line2a}).
This case is reduced to \textbf{{Case 2}}

\textbf{{Case 3.3:}} When $t\neq0$ and $1-\lambda t\neq0$, the matrix $A$
yields the linear combination of (\ref{line2}) and (\ref{line3}), due to the
following equations.%

\begin{align}
\Phi(I-A)=\left[
\begin{array}
[c]{cc}%
\varphi_{1} & \varphi_{2}\\
0 & 0
\end{array}
\right]  \left[
\begin{array}
[c]{cc}%
\lambda t & -\lambda\left(  1-\lambda t\right) \\
-t & 1-\lambda t
\end{array}
\right]  =\left[
\begin{array}
[c]{cc}%
-t\left(  \varphi_{2}-\lambda\varphi_{1}\right)  & \left(  1-\lambda t\right)
\left(  \varphi_{2}-\lambda\varphi_{1}\right) \\
0 & 0
\end{array}
\right]  ,\nonumber
\end{align}
%

\begin{align}
A\Psi= \left[
\begin{array}
[c]{cc}%
1-\lambda t & \lambda\left(  1-\lambda t\right) \\
t & \lambda t
\end{array}
\right]  \left[
\begin{array}
[c]{cc}%
\psi_{1} & 0\\
\psi_{2} & 0
\end{array}
\right]  =\left[
\begin{array}
[c]{cc}%
\left(  1-\lambda t\right)  \left(  \psi_{1}+\lambda\psi_{2}\right)  & 0\\
t\left(  \psi_{1}+\lambda\psi_{2}\right)  & 0
\end{array}
\right]  ,\nonumber
\end{align}
%

\begin{align}
\Phi A\Psi=\left[
\begin{array}
[c]{cc}%
\varphi_{1} & \varphi_{2}\\
0 & 0
\end{array}
\right]  \left[
\begin{array}
[c]{cc}%
\left(  1-\lambda t\right)  \left(  \psi_{1}+\lambda\psi_{2}\right)  & 0\\
t\left(  \psi_{1}+\lambda\psi_{2}\right)  & 0
\end{array}
\right]  =\left[
\begin{array}
[c]{cc}%
\left(  1-\lambda t\right)  \varphi_{1}\left(  \psi_{1}+\lambda\psi
_{2}\right)  +t\varphi_{2}\left(  \psi_{1}+\lambda\psi_{2}\right)  & 0\\
0 & 0
\end{array}
\right]  .\nonumber
\end{align}
Combing all the above results, we complete the proof of the Corollary.
\end{proof}

\end{document}
